% =============================================================================
% Σελίδα Τίτλου — AUTH / Τμήμα ΗΜΜΥ
% =============================================================================
\begin{titlepage}
  \newgeometry{left=2.5cm, right=2.5cm, top=1.5cm, bottom=2cm}

  \begin{center}
    \vspace*{0.3cm}
    \includegraphics[width=3.5cm]{figures/auth-logo.pdf}\\[0.6cm]
    {\large Αριστοτέλειο Πανεπιστήμιο Θεσσαλονίκης}\\[0.3em]
    {\normalsize Πολυτεχνική Σχολή}\\[0.3em]
    {\normalsize Τμήμα Ηλεκτρολόγων Μηχανικών \& Μηχανικών Υπολογιστών}\\[0.3em]
    {\small Τομέας Ηλεκτρονικής και Υπολογιστών}

    \vspace{2cm}
    {\Large\scshape Διπλωματική Εργασία}

    \vspace{1.2cm}
    \rule{3cm}{0.4pt}\;\textbullet\;\rule{3cm}{0.4pt}

    \vspace{1cm}
    {\LARGE\bfseries
      Βελτιστοποίηση του Μετα-Κβαντικού\\[0.3em]
      Κρυπτοσυστήματος HQC σε RISC-V\\[0.3em]
      μέσω Επεκτάσεων ISA}

    \vspace{0.8cm}
    \rule{3cm}{0.4pt}\;\textbullet\;\rule{3cm}{0.4pt}
  \end{center}

  \vfill

  \noindent
  \begin{minipage}[t]{0.45\textwidth}
    \textit{Εκπόνηση:}\\[0.3em]
    Φραϊδάκης Ιωάννης\\
    ΑΕΜ: 10736
  \end{minipage}%
  \hfill
  \begin{minipage}[t]{0.45\textwidth}
    \raggedleft
    \textit{Επίβλεψη:}\\[0.3em]
    Παπαευσταθίου Ιωάννης\\
    Καθηγητής Α.Π.Θ.
  \end{minipage}

  \vspace{2cm}
  \begin{center}
    \includegraphics[width=5cm]{figures/ece-logo.png}\\[0.8em]
    {\large Θεσσαλονίκη, Ιούνιος 2026}
  \end{center}

  \restoregeometry
\end{titlepage}
